Este trabalho apresentou o desenvolvimento e análise de dois sistemas fuzzy: um Mamdani para controle de ventilação e um Takagi-Sugeno para aproximação de funções não lineares~\cite{mamdani1975experiment,takagi1985fuzzy}. Ambos foram implementados do zero em Python, sem uso de bibliotecas específicas de lógica fuzzy.

O sistema Mamdani demonstrou ser eficaz para decisões qualitativas em ambientes fechados, sendo flexível e interpretável~\cite{ross2010fuzzy}. O sistema Takagi-Sugeno mostrou excelente desempenho na aproximação de funções, especialmente após otimização dos consequentes~\cite{haykin2008neural}.

A análise comparativa evidenciou a importância da escolha das funções de pertinência, operadores e técnicas de otimização para o desempenho dos sistemas fuzzy. O estudo contribui para o entendimento prático e teórico da modelagem fuzzy e pode ser expandido para aplicações mais complexas.

Como trabalhos futuros, sugere-se:
\begin{itemize}
    \item Aplicar os sistemas fuzzy em dados reais de ambientes monitorados.
    \item Explorar sistemas fuzzy com múltiplas entradas e saídas~\cite{ross2010fuzzy}.
    \item Investigar técnicas avançadas de otimização dos parâmetros fuzzy~\cite{haykin2008neural}.
    \item Integrar os sistemas fuzzy a plataformas de automação residencial ou industrial.
\end{itemize}
